\section{Signed-inverse families at prime-power steps}
\label{sec:signed-inverse}

For every sufficiently large prime power $Q=p^e$ we construct a finite family $\mathcal A_Q\subset\mathcal C$ of two-point interval components.  Its image in the simple quotient $S_Q$ is large and uniformly distributed, while its members have uniformly controlled incompatibility and reciprocal mass.

Fix the integer $\Lambda$ from Proposition~\ref{prop:prime-supply-main} and put
\[
B_{\mathrm{car}}=\Lambda+1.
\]
Let $Q=p^e$ and define
\begin{equation}\label{eq:carrier-base}
X_Q=\left\lfloor\frac{Q}{B_{\mathrm{car}}}\right\rfloor.
\end{equation}
Choose carrier primes
\begin{equation}\label{eq:carrier-band}
X_Q<b\le\Lambda X_Q,
\qquad (b,Q)=1.
\end{equation}
Writing
\[
Q=B_{\mathrm{car}}X_Q+r,
\qquad 0\le r<B_{\mathrm{car}},
\]
shows that every carrier satisfies
\begin{equation}\label{eq:carrier-exact}
B_{\mathrm{car}}b>Q,
\qquad
b\le\Lambda X_Q<B_{\mathrm{car}}X_Q\le Q.
\end{equation}
For $X_Q\ge1$ one also has
\[
B_{\mathrm{car}}X_Q\le Q<2B_{\mathrm{car}}X_Q,
\]
so $X_Q\asymp Q$.  Proposition~\ref{prop:prime-supply-main}, applied at $X_Q$, therefore gives $\asymp Q/\log Q$ possible carriers; imposing $(b,Q)=1$ removes at most the prime $p$.

For each carrier $b$, let $c_+\in\{1,\ldots,Q-1\}$ be the inverse of $b$ modulo $Q$, put $c_-=Q-c_+$, and write
\begin{equation}\label{eq:signed-inverse}
bc_+=Qk_++1,
\qquad
bc_-=Qk_--1.
\end{equation}
Then
\[
1\le k_\pm<b<Q,
\qquad
k_++k_-=b.
\]
Since $p\nmid b$, the two coefficients cannot both be divisible by $p$.  Choose a sign deterministically: if $e=1$, take the larger coefficient; if $e\ge2$, take $+$ unless $p\mid k_+$, and take $-$ otherwise.  With $k=k(Q,b)$ the chosen coefficient, define
\begin{equation}\label{eq:signed-block}
I(Q,b)=
\begin{cases}
[Qk,Qk+1],&+,\\
[Qk-1,Qk],&-.
\end{cases}
\end{equation}
If $Q=p$ is prime, this rule gives
\begin{equation}\label{eq:atomic-k-lower}
k\ge b/2>X_Q/2\gg p.
\end{equation}

\subsection{Transversality and deterministic deduplication}

The distinguished denominator $Qk$ has largest prime-power stage exactly $Q$: the coefficient $k$ is a $p$-unit, and every prime-power factor contributed by $k$ has rank below $Q$ because $k<Q$.

For the companion denominator $bc_\pm=Qk\pm1$, both factors $b$ and $c_\pm$ are below $Q$.  Hence a new prime-power factor of rank at least $Q$ can occur only if $b\mid c_\pm$.  Equivalently $b^2\mid Qk\pm1$.  Write $c_\pm=\ell b$.  From $c_\pm<Q<B_{\mathrm{car}}b$ one obtains $1\le\ell<B_{\mathrm{car}}$, and \eqref{eq:signed-inverse} gives
\begin{equation}\label{eq:bad-carrier}
\ell b^2\equiv\pm1\pmod Q.
\end{equation}
For fixed $(\ell,\pm)$ this congruence has uniformly boundedly many unit roots modulo $p^e$: at most two when $p$ is odd and at most four when $p=2$.  Thus only $O_{B_{\mathrm{car}}}(1)$ carriers fail transversality.

Different carriers can produce the same coefficient $k$.  Every carrier producing that coefficient divides one of $Qk-1,Qk+1$, both below $Q^2$.  If $Q>B_{\mathrm{car}}^3$, no one of these integers can contain three distinct carrier primes, because each carrier exceeds $Q/B_{\mathrm{car}}$.  Thus the multiplicity of a coefficient is uniformly bounded.  Keep, for each coefficient, the occurrence with least carrier.  The resulting coefficient set still has cardinality
\[
\asymp Q/\log Q.
\]

If $e\ge2$, order the remaining coefficients $k_1<\cdots<k_N$ and retain the upper half.  Since $k_i\ge i$, every retained coefficient then satisfies
\begin{equation}\label{eq:coef-lower}
k\gg Q/\log Q.
\end{equation}
If $Q=p$ is prime, retain all coefficients.  Let $\mathcal A_Q$ be the resulting family of intervals \eqref{eq:signed-block}.

\subsection{The simple-quotient value}

Modulo $F_{<Q}$, the companion reciprocal has only smaller prime-power stages, so the class contributed by $I(Q,b)$ is the class of $1/(Qk)$.  Choose $u\in\Z$ with
\[
uk\equiv1\pmod p,
\qquad uk=1+pt.
\]
Then
\begin{equation}\label{eq:simple-coordinate-calculation}
\frac{u}{Q}-\frac1{Qk}
=\frac{t}{p^{e-1}k}.
\end{equation}
Every prime-power factor of the denominator on the right has rank below $Q$: the $p$-part is at most $p^{e-1}$, and every factor coming from $k$ is below $Q$.  Hence the difference belongs to $F_{<Q}$.  Relative to the canonical generator represented by $1/Q$, the simple-quotient map
\[
\phi_Q:\mathcal A_Q\longrightarrow S_Q^\times
\]
is therefore
\begin{equation}\label{eq:simple-value}
\phi_Q(I(Q,b))=k^{-1}\pmod p.
\end{equation}
The interval $I(Q,b)$ itself recovers $Q$ and $k$: exactly one of its two denominators has largest prime-power stage $Q$, namely $Qk$.  Thus the distinguished-point map $I(Q,b)\mapsto Qk$ is globally injective.

For compatibility estimates, by an \emph{interval envelope} we mean a finite interval $J\subset\N_{>0}$ together with the one-step exclusion zone required for nonadjacency.  A member of $\mathcal A_Q$ is said to meet the envelope if its support intersects this exclusion zone.

\begin{proposition}[Uniform estimates for the signed-inverse families]\label{prop:signed-inverse-capabilities}
There is $Q_0$ such that for every prime power $Q=p^e\ge Q_0$ the following hold.
\begin{enumerate}
\item There is $\kappa>0$, independent of $Q$, for which
\begin{equation}\label{eq:family-size}
|\mathcal A_Q|\ge \kappa\frac Q{\log Q}.
\end{equation}
\item For every $s\in S_Q^\times$,
\begin{equation}\label{eq:fibre-bound}
|\phi_Q^{-1}(s)|\le Q/p.
\end{equation}
\item The incompatibility graph on $\bigcup_{Q\ge Q_0}\mathcal A_Q$ has uniformly bounded degree.  More generally, for every fixed $L\ge1$ there is $C_L$ such that an interval envelope of cardinality at most $L$ meets at most $C_L$ members of the union.
\item Every support point $z$ of a member of $\mathcal A_Q$ satisfies
\begin{equation}\label{eq:support-window}
c\frac{Q^2}{\log Q}\le z\le2Q^2
\end{equation}
for a fixed $c>0$.  If $Q$ is prime then $z\gg Q^2$.
\item There are nonnegative numbers $\omega_Q$ such that, whenever $O\subseteq\mathcal A_Q$ has $|O|\le p$ and its literal union belongs to $\mathcal C$,
\[
W\!\left(\bigcup_{I\in O}I\right)\le\omega_Q,
\]
and
\begin{equation}\label{eq:role-summable}
\sum_Q\omega_Q<\infty.
\end{equation}
\end{enumerate}
\end{proposition}

\begin{proof}
The first assertion follows from the carrier count, the bounded exceptional set, bounded coefficient multiplicity and, for $e\ge2$, the upper-half thinning.  For \eqref{eq:fibre-bound}, equality of simple-quotient values implies equality of $k$ modulo $p$.  A fixed nonzero residue class modulo $p$ occurs exactly $Q/p$ times among $1\le k<Q=p^e$, and the deterministic deduplication retains at most one interval for each coefficient.

For incompatibility, each interval occupies two adjacent integers around its globally injective distinguished point $Qk$.  If an envelope has cardinality at most $L$, every member meeting it has distinguished point in a fixed one-step enlargement of the envelope.  This enlargement contains $O(L)$ integer sites and global injectivity leaves at most one retained interval per site.

The support bounds follow from $z=Qk+O(1)$ together with \eqref{eq:coef-lower}; the prime lower bound follows from \eqref{eq:atomic-k-lower}.  Finally, a prime family has $O(p/\log p)$ intervals, each of reciprocal mass $O(p^{-2})$, so the total mass of the whole family is $O(1/(p\log p))$; this may be used for $\omega_p$.  For $Q=p^e$, $e\ge2$, each interval has mass $O(\log Q/Q^2)$, so every compatible union of at most $p$ members has mass at most
\[
\omega_{p^e}
=O\!\left(\frac{p\log Q}{Q^2}\right)
=O\!\left(\frac{e\log p}{p^{2e-1}}\right).
\]
The prime terms are summable by grouping primes into the fixed geometric bands of Proposition~\ref{prop:prime-supply-main}; the composite terms are bounded by a constant multiple of $\sum_p p^{-2}$ after summing over $e\ge2$.  Hence \eqref{eq:role-summable} holds.
\end{proof}

In particular, for every $\varepsilon>0$ there is $Q_1$ such that every finite set of prime-power steps of rank at least $Q_1$ has total reciprocal-mass bound below $\varepsilon$.
